\begin{beginnerstatement}[Kurz erklärt: Kundenanforderungen]

Eine Kundenanforderung beschreibt, was ein Produkt leisten oder unter welchen
Bedingungen es arbeiten soll.

Eine funktionale Anforderung benennt eine gewünschte Funktion; eine technische
Anforderung legt beispielsweise Plattform, Sicherheit oder Betrieb fest.

Offline-Betrieb bedeutet, dass benötigte Aufgaben auch ohne Netzwerkverbindung
möglich sein sollen.

Eine externe API ist eine Schnittstelle zu einem anderen System.

Monitoring und Observability helfen dabei, den Zustand und das Verhalten einer
laufenden Anwendung zu beobachten und Fehler zu verstehen.

Ein Backup ist eine gesicherte Kopie, aus der Daten wiederhergestellt werden
können.

Compliance-Anforderungen halten verbindliche Vorgaben fest; ein Audit ist eine
nachvollziehbare Prüfung ihrer Einhaltung.

Diese Anforderungen werden vor der Profilwahl gesammelt, weil das Template
nicht alle davon bereits abdeckt.

\end{beginnerstatement}
